\subsection{A small nilpotent leaf beside a trivial leaf}

\begin{proposition}[Proved: nilpotent four-state leaves through center eleven]
\label{CP2:A-star-nilpotent-four-leaf-eleven}
Let a source star have center $M$ of order at most eleven,
one trivial leaf on an arbitrary nonempty right act $S$,
and a leaf $N=N_{\mathrm{nil}}$ on a nonempty biset $T$ of
size at most four. It cannot have the rational homology
and cohomology ring of $\CP^2$.
\end{proposition}
\begin{proof}
Use the notation and the complete natural right-projective
$N$ resolution of
Proposition~\ref{CP2:A-nilpotent-four-corner-four-states}:
\[
 P_iP_j=P_j,\quad FP_i=P_i,\quad P_iF=P_0,
 \quad F^2=P_0,
\]
and put $U=P_0T=P_1T$, $D=(1-P_0)\mathbb QT$, and
$H=F-P_0$. The set $U$ is $M$ invariant and $N$ acts
trivially on it. If $|U|=1$, its state is fixed by both
corners. Its bridge section forces every global second
class to restrict trivially to $BM$. The leaf corners
have zero positive rational cohomology, so the two-block
relative-product argument kills all second products.
Assume $|U|\geq2$, whence $\dim D\leq2$.

Collapse the rationally acyclic corner $BN$ to a point,
then identify the two leaf points. The second operation
adds a circle to the homotopy type. The resulting actual
cofiber $Y$ therefore has bounded rational homology and
the same cohomology ring in degrees at least two as the
star. Write $R=\mathbb QM$ and
\[
 L=\ker\bigl(\mathbb QS\oplus\mathbb QU
                         \xrightarrow{\varepsilon+\varepsilon}\mathbb Q\bigr).
\]
The natural augmented bridge calculation in the cited
proposition computes its reduced chains by
$(K\otimes_R^{\mathbf L}\mathbb Q)[1]$, where
\[
 K_0=L,\quad K_j=D\ (j\geq1),\qquad
 d_1=(0,P_1-P_0),\quad d_j=H\ (j\geq2).                    \tag{1}
\]
This is the same actual bridge augmentation, with the
degree-zero summand $\mathbb QS$ added. The natural
inclusion of the sub-biset $U$ is the inclusion $L\to K$.

If $H=0$, (1) is the direct sum of $[D\to L]$ and the
complexes $D[j]$, $j\geq2$. Boundedness of $Y$ therefore
forces every $\operatorname{Tor}^R_i(D,\mathbb Q)$ to be
zero. Consequently the inclusion $L\to K$ becomes a
quasi-isomorphism after derived tensoring. On the original
stars this is the actual inclusion replacing $T$ by $U$;
the common extra circle and the collapsed acyclic corner
do not alter that homology comparison. It gives a rational
ring equivalence. Replacing the trivial $N$ action on $U$
by a trivial corner gives another actual rational ring
equivalence. Identifying its two leaves now gives the
ordinary cone on $S\sqcup U$, with the same positive
products and bounded rational homology. The bounded-eleven
act theorem makes all those products zero.

If $H\ne0$, then $\dim D=2$, $|U|=2$, and
$\ker H=\operatorname{im}H$ is one-dimensional. Consider
the full edge $X_T$ on $M,N,T$ separately. Its reduced
augmented right-$R$ complex is quasi-isomorphic to
\[
 [D/\operatorname{im}H\xrightarrow{\overline d_1}
           \ker(\mathbb QU\to\mathbb Q)],                  \tag{2}
\]
with one-dimensional terms. If $\overline d_1\ne0$, (2)
is acyclic as an $R$ complex. Thus $X_T$ is rationally
acyclic. Collapsing this actual connected subspace in
the whole star is a rational homology and ring equivalence
to the unreduced suspension of $B(S,M)$. All positive
products vanish.

Finally, if $\overline d_1=0$, then $P_1=P_0$ as functions.
The map $P_0:T\to U$ commutes with every $N$ and $M$
actor. Extend it by the identity on $S$. This gives an
actual cut retraction of the whole star onto the one with
$T$ replaced by $U$. Its relative quotient is a suspension.
The retract argument of
Lemma~\ref{CP2:A-star-general-leaf-merge} therefore retains
the rational projective-plane homology and ring on its
image: a zero second group would force all products to
be relative and zero; a nonzero second group is a split
rank-one summand whose nonzero square forces the fourth
summand as well. Collapse the now trivial $N$ action and
identify the two leaves, as above. The bounded-eleven act
theorem again kills all second products. These cases
exhaust $H^2=0$ on the vector space $D$ of dimension at
most two and prove the assertion.
\end{proof}

\begin{corollary}[Proved: a small four-corner leaf at a center through eight]
\label{CP2:A-star-eight-four-state-four-leaf}
A source star with center order at most eight, one trivial
leaf on any nonempty stalk, and another leaf of order at
most four on at most four states cannot have the rational
homology and cohomology ring of $\CP^2$.
In particular the nineteen-arrow profile $(8,1,4;2,4)$
is excluded. The opposite statements hold as well.
\end{corollary}
\begin{proof}
For a nonnilpotent leaf use
Corollary~\ref{CP2:A-star-eight-nonnilpotent-four-leaves}.
For the nilpotent leaf use the preceding proposition.
\end{proof}
